2.1 线性化与对象模型
温度对象在完整工作范围内并非线性系统：散热与温差有关，PWM 电压与实际热功率并不始终成比例，而且加热器只能升温，无法主动冷却。因此控制器设计采用工作点附近的线性增量模型。绝对温度由工作点温度和温度偏差组成，这也是 Simulink 模型中加入 25 °C 偏置的原因。
G_PT1(s) = K / (T s + 1)	(1)
G_PT2(s) = ω_n² / (s² + 2 D ω_n s + ω_n²)	(2)
2.2 由阶跃响应确定参数
对于输入幅值为 Δu 的阶跃，静态增益为 K=Δy∞/Δu。PT1 环节的时间常数 T 为从阶跃开始到输出达到总变化量 63.2% 所需的时间。PT2 环节还需要阻尼比 D 和固有角频率 ωn，可以分别由超调量和调节时间估算。分析时还必须区分对象本身的时间常数、纯延迟和传感器滞后。
M_p = exp(-π D / √(1-D²))	(3)
t_{ε,3%} ≈ 3 / (D ω_n)	(4)
2.3 代数补偿控制器
标准负反馈系统的闭环传递函数为 Gcl=KG/(1+KG)。给定期望闭环模型 GM=BM/AM，并将对象写为 G=B/A，可解得补偿控制器：
K(s) = B_M(s) A(s) / ((A_M(s) - B_M(s)) B(s))	(5)
控制器必须因果、可实现。精确零极点相消只适用于稳定且建模准确的动态。为了对阶跃设定值实现零稳态误差并抑制恒定扰动，控制器需要积分作用。
